\section[H. 一般攻击魔法]{\myboxviii{      \BF{\TT{H}}} 一般攻击魔法}

\textit{运行时间限制}   X second

\textit{内存使用限制} X Gigabyte

\textit{读入文件} \qquad 标准输入

\textit{输出文件} \qquad 标准输出


\begin{quote}
	\textit{"Zoltraak."}
\end{quote}

\subsection{题目描述}

在一个 $n \times m$ 的二维网格图中，网格中的每个格子可以处于以下三种状态之一：

\begin{enumerate}
	\item \TT{/}（正斜杠镜子）：双面反光。光线在此发生 $90^\circ$ 偏折，即连通（上方与右方）以及（下方与左方）。
	\item \TT{\char`\\}（反斜杠镜子）：双面反光。光线在此发生 $90^\circ$ 偏折，即连通（上方与左方）以及（下方与右方）。
	\item \TT{.}（空格子）：无障碍物。光线在此沿直线穿透，不发生偏折，即连通（上方与下方）以及（左方与右方）。
\end{enumerate}

初始时，网格中的部分格子已经固定放入了 \TT{/} 或 \TT{\char`\\} 镜子。其余格子目前为空。

对于每一个初始为空的格子，你可以选择\BF{将其保留为空}，或者将其\BF{替换为 \TT{/} 镜子}，或者将其\BF{替换为 \TT{\char`\\} 镜子}。

网格的最外围边界共有 $2(n + m)$ 个"端口"，光线可以从这些边界端口射入或射出网格。

现给定一束光线的\BF{初始射入端口}（及射入方向）以及期望的\BF{目标射出端口}（及射出方向）。

请问是否存在一种合法的网格填充方案，使得光线从指定的射入端口进入后，在网格内经过若干次直线传播与镜面反射，最终\BF{恰好}从指定的目标射出端口射出？

如果存在合法方案，请输出填充后的网格；如果存在多种方案，输出任意一种即可。如果不可能做到，请报告无解。

\subsection{输入格式}

第一行输入一个整数 $T$ ($1 \le T \le 100$)，表示测试数据的数量。

每组测试数据第一行输入两个整数 $n$ 和 $m$ ($1 \le n \times m \le 2 \times 10^5$)，表示网格的行数和列数。

随后 $n$ 行，每行输入一个长度为 $m$ 的字符串，表示网格的初始状态。其中 \TT{/} 表示正斜杠镜子，\TT{\char`\\} 表示反斜杠镜子，\TT{.} 表示空格子。

随后一行输入六个整数 $x_1, y_1, d_1, x_2, y_2, d_2$，其中：

\begin{itemize}
	\item $(x_1, y_1)$ 表示光线射入端口的位置，$d_1 \in \{0, 1, 2, 3\}$ 表示射入方向（$0$ 表示向右，$1$ 表示向下，$2$ 表示向左，$3$ 表示向上）。
	\item $(x_2, y_2)$ 表示光线期望射出端口的位置，$d_2 \in \{0, 1, 2, 3\}$ 表示射出方向（定义同上）。
\end{itemize}

端口坐标的定义为：$(0, j)$ 表示第 $j$ 列的上边界，$(n + 1, j)$ 表示第 $j$ 列的下边界，$(i, 0)$ 表示第 $i$ 行的左边界，$(i, m + 1)$ 表示第 $i$ 行的右边界。保证输入的端口在边界上，且射入方向指向网格内部，射出方向指向网格外部。

保证所有测试数据的 $\sum (n \times m) \le 2 \times 10^5$ 。

\subsection{输出格式}

对于每组测试数据,若存在合法方案，输出 $n$ 行，每行输出一个长度为 $m$ 的字符串，表示填充后的网格。否则，输出一行 \TT{IMPOSSIBLE}。


\subsection{样例 1}

\begin{minipage}{\linewidth}
	\begin{minipage}[t]{0.49\linewidth}
		\begin{lstlisting}[title=\texttt{standard input}]
1
2 2
..
..
1 0 0 0 2 3
\end{lstlisting}
	\end{minipage}
	\begin{minipage}[t]{0.49\linewidth}
		\begin{lstlisting}[title=\texttt{standard output}]
./
\.
\end{lstlisting}
	\end{minipage}
\end{minipage}


\subsection{样例 1 解释}

对于第 $1$ 组测试数据，光线从第 $1$ 行的左边界射入（向右），经过 $(1, 1)$ 遇到 \TT{.} 继续向右，经过 $(1, 2)$ 遇到 \TT{/} 向上偏折，从第 $2$ 列的上边界射出（向上）。
